% 1일차 실습 4 — 임베딩 → SQL로 묻기 → 더 복잡한 질의 (실습_SQL_임베딩_질의.md)
\section{실습 4. 임베딩 → SQL로 묻기}
% 글꼴에 ⓪(U+24EA)이 없어 TikZ로 대신 그린다
\providecommand{\circzero}{\tikz[baseline=(z.base)]\node[circle,draw,line width=0.4pt,inner sep=0.4pt,minimum size=0.95em,font=\scriptsize](z){0};}

% ─────────────────────────────────────────────
\begin{frame}{임베딩 → SQL로 묻기 → 더 복잡한 질의 \src{\normalsize 1일차 마무리}}
\begin{block}{목표}
글을 숫자로 바꾸는 \textbf{임베딩}이 무엇인지 손으로 확인하고, JPMorgan 상품 담당자가 되어 상품 라인업을 \textbf{SQL로} 분석한 뒤, \textbf{벡터로만} 할 수 있는 분석(특이 상품 찾기, 새 상품과 닮은 상품 찾기)을 해 본다.
\end{block}
\vspace{0.3em}
\footnotesize
\begin{tabularx}{\textwidth}{@{}l X r@{}}
\toprule
\textbf{순서} & \textbf{내용} & \textbf{시간} \\
\midrule
\tagbox{accent}{A} & Azure OpenAI 만들고 \textbf{벡터화 시작} (약 19분 걸리므로 먼저 시작해 두고 B·C 동안 돌린다) & 8분 \\
\tagbox{blue}{B} & 임베딩이 무엇인가: TensorFlow Embedding Projector로 보기, 단어 계산, 문장 유사도 & 15분 \\
\tagbox{teal}{C} & 상품군 분석: JPMorgan 상품 담당자로서 ``무엇을 얼마나 팔았나'' (SQL, 벡터 불필요) & 12분 \\
\tagbox{ink!55}{D} & 벡터로만 할 수 있는 분석: 키워드 vs 뜻(3분), 특이 상품 찾기, 새 상품 설계안과 닮은 상품 찾기 & 18분 \\
\tagbox{muted}{E} & 1일차 정리 & 5분 \\
\bottomrule
\end{tabularx}
\end{frame}

% ═════════════════════════════════════════════ A
\begin{frame}[fragile]{\tagbox{accent}{A} Azure OpenAI 만들기}
VS Code 터미널 \src{(리소스 그룹 \texttt{jpm-course}, 이름은 \texttt{jpm-aoai-<이니셜>-<숫자>})}
\vspace{0.4em}
\begin{Verbatim}
az cognitiveservices account create -g jpm-course -n jpm-aoai-<이니셜>-<숫자> -l japaneast \
  --kind OpenAI --sku S0 --custom-domain jpm-aoai-<이니셜>-<숫자> --yes
az cognitiveservices account deployment create -g jpm-course -n jpm-aoai-<이니셜>-<숫자> \
  --deployment-name emb3s --model-name text-embedding-3-small --model-version 1 --model-format OpenAI \
  --sku-name GlobalStandard --sku-capacity 1000
\end{Verbatim}
\vspace{0.6em}
\begin{columns}[T]
\begin{column}{0.48\textwidth}
\begin{block}{만드는 것}
\small 계정 하나 + 임베딩 모델 배포 \texttt{emb3s}\\
\src{모델 \texttt{text-embedding-3-small}, 지역 \texttt{japaneast}}
\end{block}
\end{column}
\begin{column}{0.48\textwidth}
\begin{alertblock}{먼저 시작해 둔다}
\small 벡터화는 \num{약 19분}. B·C 동안 돌린다.
\end{alertblock}
\end{column}
\end{columns}
\end{frame}

\begin{frame}[fragile]{\tagbox{accent}{A} \texttt{InsufficientQuota} 오류가 나면}
\begin{alertblock}{무료 계정은 Azure OpenAI를 \textbf{하나만}}
지운 계정도 \num{48시간} 동안 ``삭제 대기''로 남아 그 자리를 차지한다\\
\src{(예: 연습으로 만들었다 지운 계정)}
\end{alertblock}
\vspace{0.4em}
휴지통을 확인하고 완전히 지운 뒤 앞의 명령을 다시 실행한다.
\vspace{0.4em}
\begin{Verbatim}
az cognitiveservices account list-deleted -o table
az cognitiveservices account purge -l japaneast -g <그 계정의 리소스 그룹> -n <그 계정 이름>
\end{Verbatim}
\end{frame}

\begin{frame}[fragile]{\tagbox{accent}{A} 키 적고 벡터화 시작}
\small
\texttt{.env}에 \textbf{직접} 적는다.
\begin{itemize}\setlength\itemsep{1pt}
\item \texttt{AOAI\_ENDPOINT=https://jpm-aoai-<이니셜>-<숫자>.openai.azure.com/}
\item \texttt{AOAI\_KEY} \src{(포털 → 이 리소스 → \textbf{키 및 엔드포인트} → 키 1)}
\end{itemize}
\vspace{0.3em}
\begin{Verbatim}
uv run jpm_kit/embed.py --try 3
\end{Verbatim}
청크 3개가 각각 \num{숫자 1,536개}로 바뀌는 것을 본다. 그다음 Claude Code 패널:
\vspace{0.3em}
\begin{Verbatim}
jpm_kit/embed.py 를 백그라운드로 실행하고, jpm_kit/out/embed_progress.txt 로 진행 상황을 알려줘.
\end{Verbatim}
\end{frame}

% ═════════════════════════════════════════════ B
\begin{frame}{\tagbox{blue}{B} 임베딩이 무엇인가 \src{\normalsize 벡터화가 도는 동안}}
\begin{block}{한 줄 정의}
글(단어·문장·문단)을 \textbf{숫자 목록(벡터)}으로 바꾸는 것. 뜻이 비슷한 글은 비슷한 숫자 목록이 되도록 학습되어 있어서, \textbf{숫자끼리의 거리로 뜻이 얼마나 가까운지} 잴 수 있다.
\end{block}
\vspace{0.5em}
\centering
\begin{tikzpicture}
\node[box=ink,minimum width=3.2cm] (t1) {글 하나};
\node[box=ink,minimum width=3.2cm,below=0.5cm of t1] (t2) {뜻이 비슷한 글};
\node[box=blue,minimum width=5.6cm,right=1.8cm of t1,font=\small\ttfamily] (v1) {[ 숫자, 숫자, 숫자, … ]};
\node[box=blue,minimum width=5.6cm,right=1.8cm of t2,font=\small\ttfamily] (v2) {[ 숫자, 숫자, 숫자, … ]};
\draw[arr] (t1) -- node[above,font=\scriptsize]{임베딩} (v1);
\draw[arr] (t2) -- node[above,font=\scriptsize]{임베딩} (v2);
\draw[accent,line width=0.9pt,{Stealth[length=5pt]}-{Stealth[length=5pt]}] ([xshift=2mm]v1.east) to[bend left=50] node[right,font=\scriptsize\bfseries,accent,align=left]{거리가\\가깝다} ([xshift=2mm]v2.east);
\end{tikzpicture}
\end{frame}

\begin{frame}{\tagbox{blue}{B} 말 네 개}
\begin{tabularx}{\textwidth}{@{}l X@{}}
\toprule
\textbf{말} & \textbf{뜻} \\
\midrule
벡터 & 숫자 목록. 단어 벡터는 \num{200개}, 오늘 쓰는 문장 임베딩은 \num{1,536개} \\
차원 & 숫자의 개수. 숫자 하나하나에는 사람이 읽을 뜻이 없다 \\
코사인 유사도 & 두 벡터의 방향이 얼마나 같은가. \textbf{1}이면 같은 방향(뜻이 가까움), \textbf{0}이면 무관 \\
벡터 검색 & 질문도 같은 모델로 벡터로 바꾼 뒤, 가장 가까운 청크를 찾는 것 \\
\bottomrule
\end{tabularx}
\end{frame}

\begin{frame}{\tagbox{blue}{B} \circzero{} 눈으로 보기: TensorFlow Embedding Projector \src{\normalsize 브라우저, 5분}}
\small
구글 TensorFlow 팀이 공개한 임베딩 시각화 도구. 위키백과 글로 학습한 단어 \textbf{1만 개}의 벡터(단어마다 숫자 200개)를 3차원으로 눌러 점으로 보여 준다.
\begin{enumerate}\setlength\itemsep{3pt}
\item \texttt{projector.tensorflow.org}를 연다. 왼쪽 위 데이터가 \textbf{Word2Vec 10K}인지 확인한다(기본값).
\item 오른쪽 \textbf{Search} 칸에 \texttt{bank} → 아래 목록에서 \textbf{bank를 클릭} → 오른쪽에 \textbf{Nearest points in the original space}(숫자 200개 공간에서 가장 가까운 단어)와 거리가 나온다.\\
\src{이 숫자는 \textbf{코사인 거리}라 \textbf{작을수록 가깝다} (① 스크립트가 보여 주는 \textbf{유사도}는 클수록 가까움).}
\item \textbf{Isolate 101 points}를 눌러 이웃만 남기고, 점 구름을 마우스로 돌려 본다. \texttt{dollar}, \texttt{bond}, \texttt{risk}도 찾아본다.
\item 왼쪽 아래 투영 방법을 \textbf{PCA → T-SNE}로 바꿔, 비슷한 단어끼리 무리 짓는 모습을 본다.
\end{enumerate}
\end{frame}

\begin{frame}{\tagbox{blue}{B} \circzero{} 적을 것}
\begin{tabularx}{\textwidth}{@{}X p{6cm}@{}}
\toprule
\textbf{확인} & \textbf{내 값} \\
\midrule
\texttt{bank}의 가장 가까운 단어 5개 & \\[4pt]
\texttt{bank} 이웃 중 금융이 아닌 단어 (어떤 뜻의 bank?) & \\[4pt]
\texttt{bond}의 가장 가까운 단어 3개 (채권 뜻인가?) & \\[4pt]
\bottomrule
\end{tabularx}
\vspace{0.6em}
\begin{alertblock}{단어 하나에 숫자 묶음이 하나뿐}
\small \texttt{bank} 이웃에는 banks·monetary·imf 같은 금융 단어 사이에 river·thames·canal(강둑), gaza·west·strip(요르단강 서안 West Bank)이 섞여 나온다. 뜻이 섞인 것이다 → ②에서 문장 임베딩이 이것을 어떻게 푸는지 본다.
\end{alertblock}
\end{frame}

\begin{frame}[fragile]{\tagbox{blue}{B} ① 같은 자료로 계산하기: 단어 계산}
\small
AI 호출 없음, 내 노트북에서 numpy로. \circzero{}에서 본 \textbf{바로 그 Word2Vec 10K 자료}(TensorFlow Embedding Projector 공개 파일, 약 8MB)를 내려받아 직접 계산한다. 핵심은 아래 몇 줄.
\vspace{0.3em}
\begin{Verbatim}
idx = {word: i for i, word in enumerate(words)}          # 단어 → 행 번호
q = V[idx["king"]] - V[idx["man"]] + V[idx["woman"]]     # 벡터 산술
scores = (V @ q) / (np.linalg.norm(V, axis=1) * np.linalg.norm(q))   # 모든 단어와 코사인 유사도
# scores 가 가장 큰 단어 = 답 (식에 쓴 세 단어는 뺀다)
\end{Verbatim}
\vspace{0.3em}
\begin{Verbatim}
uv run jpm_kit/emb_1_words.py
\end{Verbatim}
\end{frame}

\begin{frame}{\tagbox{blue}{B} ① 무엇이 나오나}
\begin{columns}[T]
\begin{column}{0.48\textwidth}
\begin{block}{가까운 단어}
\texttt{dollar → currency, usd, franc …}\\
\src{\circzero{}의 목록과 비교해 보기}
\end{block}
\begin{block}{단어 계산}
\texttt{king - man + woman = queen}\\
\texttt{paris - france + japan = tokyo}\\[2pt]
→ 관계가 \textbf{방향}으로 들어 있다
\end{block}
\end{column}
\begin{column}{0.48\textwidth}
\begin{alertblock}{한계}
\texttt{bond}의 이웃이 hydrogen, molecule(화학 결합).\\[2pt]
\texttt{bank}는 loan(\num{0.50})보다 river(\num{0.56}, 강둑)에 가깝다\\[2pt]
→ 단어 하나에 숫자 묶음이 하나뿐이라 뜻이 섞인다
\end{alertblock}
\src{\texttt{jpm\_kit/out/word\_map.png}: 200개 숫자를 2차원으로 눌러 본 지도}
\end{column}
\end{columns}
\end{frame}

\begin{frame}[fragile]{\tagbox{blue}{B} ① 왕·여왕·남자·여자 그림 \src{\normalsize \texttt{jpm\_kit/out/king\_queen.png}}}
\begin{columns}[c]
\begin{column}{0.44\textwidth}
\centering
\begin{tikzpicture}[x=1cm,y=1cm]
\draw[muted,-{Stealth[length=4pt]}] (0,0) -- (4.4,0) node[below left,font=\scriptsize,align=right]{남자→여자\\(woman - man)};
\draw[muted,-{Stealth[length=4pt]}] (0,0) -- (0,3.2) node[above,font=\scriptsize]{보통 사람→왕족};
\node[circle,fill=ink,inner sep=1.5pt,label={[font=\scriptsize]below:man}] (m) at (0.9,0.7) {};
\node[circle,fill=ink,inner sep=1.5pt,label={[font=\scriptsize]below:woman}] (w) at (3.3,0.7) {};
\node[circle,fill=accent,inner sep=1.5pt,label={[font=\scriptsize]above:king}] (k) at (0.9,2.4) {};
\node[circle,fill=accent,inner sep=1.5pt,label={[font=\scriptsize]above:queen}] (q) at (3.3,2.4) {};
\draw[arr,ink] (m) -- (w);
\draw[arr,accent] (k) -- (q);
\end{tikzpicture}\\
\src{화살표 방향만 나타낸 모식도}
\end{column}
\begin{column}{0.52\textwidth}
\small
가로축 = 남자→여자 방향(woman - man), 세로축 = 보통 사람→왕족 방향.\\[3pt]
man→woman 화살표와 \textbf{king→queen 화살표가 같은 쪽(오른쪽)}을 향한다. 그래서 king에서 man을 빼고 woman을 더하면 queen 근처에 닿는다.\\[3pt]
prince→princess, boy→girl, father→mother도 같은 쪽.\\[3pt]
\alert{같은 쪽을 향하지 않는 쌍(husband→wife)}도 찾아 왜 그럴지 생각해 본다.
\end{column}
\end{columns}
\vspace{0.4em}
\src{Projector에서 같은 그림을 보려면 (강사 시연): 왼쪽 아래 \textbf{CUSTOM} → 네 칸 오른쪽의 \texttt{.*} 버튼을 눌러 켠 뒤 Left \texttt{\^{}man\$}, Right \texttt{\^{}woman\$}, Up \texttt{\^{}(king|queen)\$}, Down \texttt{\^{}(man|woman)\$}}
\end{frame}

\begin{frame}[fragile]{\tagbox{blue}{B} ① 직접 해 보기}
\small
두 명령을 \textbf{따로} 실행한다. \alert{한 줄에 같이 쓰면 \texttt{-{}-near}만 실행된다.}
\vspace{0.3em}
\begin{Verbatim}
uv run jpm_kit/emb_1_words.py --calc tokyo japan france
uv run jpm_kit/emb_1_words.py --near loan
\end{Verbatim}
\vspace{0.4em}
\begin{tabularx}{\textwidth}{@{}l X@{}}
\toprule
\texttt{-{}-calc A B C} & A - B + C에 가장 가까운 단어 \src{(tokyo - japan + france → paris)} \\
\texttt{-{}-near 단어} & 그 단어와 가장 가까운 단어 6개 \src{(loan → loans, debts, payments …)} \\
\bottomrule
\end{tabularx}

\vspace{0.6em}
다른 영어 단어로도 해 본다.
\end{frame}

\begin{frame}[fragile]{\tagbox{blue}{B} ② 문장을 숫자로 \src{\normalsize 내 Azure OpenAI, 1원도 안 듦}}
\begin{Verbatim}
uv run jpm_kit/emb_2_sentences.py
\end{Verbatim}
\vspace{0.4em}
\begin{tabularx}{\textwidth}{@{}X p{5.5cm}@{}}
\toprule
\textbf{확인} & \textbf{내 값} \\
\midrule
같은 bond: 채권 문장끼리 vs 채권–화학 문장 & \hfill / \hfill\null \\[6pt]
한국어 질문과 가장 가까운 영어 공시 문장 & \\[4pt]
금리만 다른 두 문장(10.75\% vs 8.10\%)의 유사도 & \\[4pt]
\bottomrule
\end{tabularx}

\vspace{0.5em}
\small 내 문장으로:
\begin{Verbatim}
uv run jpm_kit/emb_2_sentences.py "원금 보장" "principal protected" "capital at risk"
\end{Verbatim}
\end{frame}

\begin{frame}{\tagbox{blue}{B} ② 문장 임베딩이 보여 주는 것}
\begin{columns}[T]
\begin{column}{0.31\textwidth}
\begin{block}{문맥}
\textbf{문맥}이 뜻을 정한다\\
\src{단어 벡터의 한계를 넘음}
\end{block}
\end{column}
\begin{column}{0.31\textwidth}
\begin{block}{언어}
\textbf{언어가 달라도} 뜻이 같으면 가깝다
\end{block}
\end{column}
\begin{column}{0.31\textwidth}
\begin{alertblock}{숫자}
대신 \textbf{정확한 숫자는 거의 구별하지 못한다}
\end{alertblock}
\end{column}
\end{columns}
\vspace{1em}
\begin{center}
→ 그래서 \num{키워드 검색과 합친다}.
\end{center}
\vspace{0.6em}
\src{벡터화는 계속 돌고 있다. C(SQL)를 마칠 무렵 끝난다. 결과는 D 전에 적는다.}
\end{frame}

% ═════════════════════════════════════════════ C
\begin{frame}{\tagbox{teal}{C} 상품군 분석: 우리는 무엇을 얼마나 팔았나 \src{\normalsize 벡터화가 도는 동안}}
\begin{block}{역할}
여러분은 \textbf{JPMorgan 구조화상품 담당자}. 공용 상품 DB(\texttt{products}, 2026년 1–10월 \num{8,868개})로 우리 상품 라인업을 본다.\\[3pt]
상품군 이름을 이미 알고 있으므로 \textbf{벡터가 필요 없다. SQL이면 된다.}
\end{block}
\vspace{0.5em}
\textbf{공용 상품 DB 연결} \src{— PostgreSQL 확장에 두 번째 연결 (교내 네트워크에서만)}\\[4pt]
\small
\begin{tabular}{@{}l l@{}}
\toprule
서버 & \texttt{jpm-shared-d32c.postgres.database.azure.com} \\
사용자 & \texttt{student\_ro} \\
비밀번호 & \texttt{20261128} \\
데이터베이스 & \texttt{jpm} \\
\bottomrule
\end{tabular}
\end{frame}

\begin{frame}{\tagbox{teal}{C-1} 내 손으로 SQL 두 개 \src{\normalsize VS Code 쿼리 창, 5분}}
\footnotesize
\begin{tabularx}{\textwidth}{@{}c >{\raggedright\arraybackslash}p{4.2cm} >{\raggedright\arraybackslash}X p{2.6cm}@{}}
\toprule
\textbf{\#} & \textbf{질문} & \textbf{힌트} & \textbf{내 답} \\
\midrule
1 & 상품군(\texttt{product\_type})별 상품 수와 비중 & \texttt{count(*)}, \texttt{round(100.0 * count(*) / sum(count(*)) over (), 1)}, \texttt{group by 1 order by 2 desc} & \\[4pt]
2 & 조건부 쿠폰 상품의 기초자산 구조(\texttt{underlier\_structure})별 평균 연 쿠폰 — \textbf{분모(쿠폰이 있는 상품 수)도 함께} & \texttt{where product\_type = 'contingent\_income' and coupon\_annual\_pct is not null group by …} & \\[4pt]
\bottomrule
\end{tabularx}
\vspace{0.8em}
\begin{block}{막히면}
\small Claude Code 패널에 \textit{``이 질문을 SQL로 어떻게 쓰는지 설명만 해줘''}라고 묻고, 받은 SQL을 \textbf{읽고 이해한 뒤} 실행한다.
\end{block}
\end{frame}

\begin{frame}[fragile]{\tagbox{teal}{C-2} Claude에게 상품군 지도 맡기기 \src{\normalsize Claude Code 패널, 약 4분}}
\begin{Verbatim}[fontsize=\footnotesize]
너는 JPMorgan 구조화상품 담당자다. 공용 상품 DB(products, 2026년 1~10월)로 "우리가 무엇을 얼마나 팔았나" 상품군 지도를 만들어줘.
1. 상품군(product_type)별 상품 수·비중·평균 만기·평균 연 쿠폰·발행가 대비 추정가치 차이(price_value_gap_pct) 평균·원금 합계. 평균마다 분모(값이 있는 상품 수)도.
2. 월별 상품 수와 상위 4개 상품군의 월별 추이.
3. 가장 많이 쓰인 기초자산(tickers) 10개.
4. 이 숫자로 말할 수 있는 것 3가지와, 이 데이터로는 말할 수 없는 것 2가지.
DB 연결은 jpm_kit/kit_db.py 의 connect("SHARED"). 짧은 코드는 jpm_kit/out/ 에 파일로 만들어 uv run 으로 실행하고, 쓴 SQL 을 모두 보여줘. .env 는 열지 마.
\end{Verbatim}
\end{frame}

\begin{frame}{\tagbox{teal}{C-2} 확인}
\begin{tabularx}{\textwidth}{@{}X p{5.6cm}@{}}
\toprule
\textbf{확인} & \\
\midrule
Claude의 상품군별 상품 수가 내 C-1의 1번과 같은가 & \\[8pt]
평균 쿠폰·숨은 비용(gap)의 \textbf{분모}가 상품 수와 다른 상품군은? 왜 다른가 & \\[8pt]
Claude가 꼽은 ``이 데이터로는 말할 수 없는 것'' 하나 & \\[8pt]
\bottomrule
\end{tabularx}
\end{frame}

% ═════════════════════════════════════════════ D
\begin{frame}{\tagbox{ink!55}{D} 벡터로만 할 수 있는 분석 \src{\normalsize 벡터화가 끝난 뒤}}
먼저 벡터화 결과를 적는다 \src{(\texttt{jpm\_kit/out/embed\_progress.txt} 마지막 줄)}
\vspace{0.6em}
\begin{tabularx}{\textwidth}{@{}X p{7cm}@{}}
\toprule
\textbf{벡터화 결과} & \textbf{내 값} \\
\midrule
청크 수 / 걸린 시간 / 토큰 / 비용 & \\[8pt]
DB 크기 (전 → 후) & \\[8pt]
\bottomrule
\end{tabularx}
\end{frame}

\begin{frame}{\tagbox{ink!55}{D} 키워드 검색 vs 벡터 검색}
\small
\begin{tabularx}{\textwidth}{@{}l X X@{}}
\toprule
 & \textbf{키워드 검색} & \textbf{벡터 검색} \\
\midrule
찾는 기준 & 질문과 \textbf{같은 단어}가 들어 있는가 \src{(\texttt{tsv} 키워드 색인)} & 질문과 \textbf{뜻}이 가까운가 \src{(\texttt{emb} 벡터의 코사인 거리)} \\
잘 찾는 것 & 공시 용어 그대로, 정확한 숫자, 고유명사 & 바꿔 말한 질문 (공시와 다른 말로 물어도) \\
놓치는 것 & 질문의 단어가 한 문단에 모두 있어야 찾음 & 숫자만 다른 문장을 구별 못 함 (B의 10.75\% vs 8.10\%), 한국어로 막연히 물으면 자주 실패 \\
\bottomrule
\end{tabularx}
\end{frame}

\begin{frame}{\tagbox{ink!55}{D} 하이브리드: 두 결과의 \textbf{순위}를 합친다}
\begin{columns}[c]
\begin{column}{0.56\textwidth}
\centering
\begin{tikzpicture}
\node[box=accent,minimum width=3cm] (k) at (0,0.85) {키워드 검색 결과\\[-1pt]\scriptsize 1위, 2위, 3위, …};
\node[box=blue,minimum width=3cm] (v) at (0,-0.85) {벡터 검색 결과\\[-1pt]\scriptsize 1위, 2위, 3위, …};
\node[box=ink,minimum width=2.6cm,minimum height=2.2cm] (h) at (4.6,0) {하이브리드\\[-1pt]\scriptsize 점수 합으로\\[-3pt]\scriptsize 다시 줄 세움};
\draw[arr] (k.east) -- node[above,sloped,font=\scriptsize]{$1/(60+r)$} (h.west);
\draw[arr] (v.east) -- node[below,sloped,font=\scriptsize]{$1/(60+r)$} (h.west);
\end{tikzpicture}
\end{column}
\begin{column}{0.40\textwidth}
\begin{block}{RRF}
\small \texttt{ask.py}의 세 번째 결과.\\[2pt]
각 목록에서 \textbf{r위}면 \num{1 / (60 + r)} 점.
\end{block}
\begin{alertblock}{늘 이기지는 않는다}
\small 양쪽에 모두 오른 문단이 위로 올라오지만, 늘 이기지는 않는다.
\end{alertblock}
\end{column}
\end{columns}
\end{frame}

\begin{frame}[fragile]{\tagbox{ink!55}{D-1} 같은 뜻, 다른 말: \texttt{ask.py}로 3분}
\begin{Verbatim}
uv run jpm_kit/ask.py "Can the bank end the investment early?" -k 3
uv run jpm_kit/ask.py "6.0% per annum daily deduction" -k 3
\end{Verbatim}
\vspace{0.6em}
\begin{tabularx}{\textwidth}{@{}l X X@{}}
\toprule
\textbf{질문} & \textbf{키워드 1위} & \textbf{벡터 1위} \\
\midrule
고객의 말 (공시에 없는 표현) & & \\[8pt]
공시 표현 그대로 & & \\[8pt]
\bottomrule
\end{tabularx}
\end{frame}

\begin{frame}{\tagbox{ink!55}{D-1} 첫 질문에서 키워드가 아무것도 못 찾는 이유}
\begin{block}{키워드 검색은 질문의 \textbf{모든 단어}가 한 문단에 있어야 한다}
\texttt{'bank' \& 'end' \& 'invest' \& 'earli'}\\[3pt]
벡터는 \textbf{``Early redemption risk''} 문단을 찾는다.
\end{block}
\vspace{0.6em}
\begin{alertblock}{벡터가 값을 하는 때}
\textbf{묻는 말과 공시의 말이 다를 때}.\\[3pt]
아래 두 분석(D-2, D-3)이 바로 그런 경우: 무엇을 찾을지 \textbf{이름을 모르는} 분석.
\end{alertblock}
\end{frame}

\begin{frame}{\tagbox{ink!55}{D} 공시 한 건 = 벡터 하나 \src{\normalsize \texttt{jpm\_kit/similar.py}}}
\centering
\begin{tikzpicture}
\node[box=ink,minimum width=2.4cm,minimum height=2.6cm] (doc) {공시 1건};
\node[box=muted,minimum width=2.6cm,minimum height=0.55cm,right=1.2cm of doc,yshift=0.85cm,font=\scriptsize] (c1) {청크 벡터};
\node[box=muted,minimum width=2.6cm,minimum height=0.55cm,right=1.2cm of doc,font=\scriptsize] (c2) {청크 벡터};
\node[box=muted,minimum width=2.6cm,minimum height=0.55cm,right=1.2cm of doc,yshift=-0.85cm,font=\scriptsize] (c3) {청크 벡터 …};
\node[box=blue,minimum width=2.4cm,right=1.2cm of c2] (avg) {평균\\[-1pt]\scriptsize = 벡터 하나};
\draw[arr] (doc.east) -- (c1.west);
\draw[arr] (doc.east) -- (c2.west);
\draw[arr] (doc.east) -- (c3.west);
\draw[arr] (c1.east) -- (avg.west);
\draw[arr] (c2.east) -- (avg.west);
\draw[arr] (c3.east) -- (avg.west);
\node[font=\scriptsize,muted,below=0.15cm of c3] {공통 문구가 아닌 청크만};
\node[box=accent,minimum width=3cm,right=0.8cm of avg,font=\small] (cmp) {코사인 유사도로\\[-1pt]비교};
\draw[arr] (avg) -- (cmp);
\end{tikzpicture}
\vspace{0.8em}
\begin{itemize}\setlength\itemsep{2pt}\small
\item 공시 \num{1,859건}을 벡터 \num{1,859개}로 만든다 \src{(처음 한 번 약 30초)}
\item 공시끼리, 또는 새 설명과 코사인 유사도로 비교한다
\end{itemize}
\end{frame}

\begin{frame}[fragile]{\tagbox{ink!55}{D-2} 우리가 판 상품 중 특이한 것이 있나 \src{\normalsize 약 5분}}
먼저 직접 실행해 본다.
\begin{Verbatim}
uv run jpm_kit/similar.py --outliers 5
\end{Verbatim}
그다음 Claude Code 패널에서:
\begin{Verbatim}[fontsize=\footnotesize]
너는 JPMorgan 구조화상품 담당자다. 10월에 우리가 낸 공시 중 특이한 상품이 있는지 보려고 한다.
1. uv run jpm_kit/similar.py --outliers 5 를 실행하고 결과를 표로 보여줘.
2. 그 5건이 어떤 상품인지 uv run jpm_kit/ask.py "<검색어>" --accession <번호> 로 원문을 찾아 확인하고, 왜 특이한지 한 줄씩 써줘.
   문서 양식(판매 채널)이 달라서인지, 상품 구조가 달라서인지 구분하고, 근거 [accession #chunk_id] 를 붙여.
.env 는 열지 마.
\end{Verbatim}
\end{frame}

\begin{frame}{\tagbox{ink!55}{D-2} 확인}
\begin{tabularx}{\textwidth}{@{}X p{5.6cm}@{}}
\toprule
\textbf{확인} & \\
\midrule
가장 닮은 다른 공시와 0.99 이상으로 닮은 공시 비율 → 무슨 뜻인가 & \\[8pt]
특이 5건 중 \textbf{상품 구조}가 특이한 것 / \textbf{문서 양식(판매 채널)}이 달라서인 것 & \hfill / \hfill\null \\[8pt]
특이 5건 중 공용 상품 표에 없거나 분류 미확정인 것은 몇 건? 왜 그럴까 & \\[8pt]
\bottomrule
\end{tabularx}
\end{frame}

\begin{frame}[fragile]{\tagbox{ink!55}{D-3} 새 상품을 설계했다: 기존 상품 중 무엇과 닮았나 \src{\normalsize 약 5분}}
\small Claude Code 패널에서:
\begin{Verbatim}[fontsize=\footnotesize]
너는 JPMorgan 구조화상품 담당자다. 새 상품을 설계했다:
"금과 비트코인 ETF를 같은 비중으로 담은 바스켓에 연동, 만기 5년, 원금 100% 보장, 바스켓 상승분의 80% 참여, 수익 상한 40%, 이자 없음."
1. 이 설명을 공시에서 쓰는 영어 문장 3~4개로 바꿔 보여줘.
2. 비교 기준으로 평범한 상품 설명(Nasdaq-100·Russell 2000 worst-of 자동조기상환 조건부 쿠폰)도 영어로 하나 써서,
   uv run jpm_kit/similar.py --text "<새 상품>" --text "<평범한 상품>" 을 실행해.
3. 새 상품과 가장 비슷한 공시 3건이 무엇 때문에 닮았고 무엇이 다른지 ask.py --accession 으로 원문을 확인해 표로. 특히 "원금 100% 보장"이 닮은 이유에 들어 있는지 봐.
4. 이 새 상품은 기존 상품의 변형인가, 새로운 구조인가? 근거와 함께 3줄로.
.env 는 열지 마.
\end{Verbatim}
\end{frame}

\begin{frame}{\tagbox{ink!55}{D-3} 확인}
\begin{tabularx}{\textwidth}{@{}X p{5.6cm}@{}}
\toprule
\textbf{확인} & \\
\midrule
새 상품 / 평범한 상품의 1위 유사도 & \hfill / \hfill\null \\[8pt]
새 상품과 가장 닮은 공시가 닮은 이유 (무엇이 같은가) & \\[8pt]
``원금 100\% 보장''이 닮은 이유에 들어갔나? 들어가지 않았다면 그것이 벡터에 대해 말해 주는 것은 & \\[8pt]
\bottomrule
\end{tabularx}
\vspace{0.6em}
\begin{alertblock}{한국어 설명은 먼저 영어 공시 문장으로}
\small 한국어 설명을 그대로 \texttt{-{}-text}에 넣으면 영어 공시와 잘 비교되지 않는다 \src{(리허설: 유사도 \num{0.35})}. 그래서 Claude에게 먼저 영어 공시 문장으로 바꾸게 한다.
\end{alertblock}
\end{frame}

% ═════════════════════════════════════════════ E
\begin{frame}{\tagbox{muted}{E} 1일차 정리 \src{\normalsize 5줄 메모}}
\begin{enumerate}\setlength\itemsep{7pt}
\item 오늘 내 손으로 만든 것: SEC 원문 \underline{\hspace{1.2cm}}건 → 텍스트 \underline{\hspace{1.2cm}}건 → 내 DB의 표 \underline{\hspace{1.2cm}}개 (행 수)
\item 변환기가 바꾸지 \textbf{않는} 것 한 가지
\item 키워드 검색이 잘 찾는 질문 / 벡터 검색이 잘 찾는 질문
\item 벡터가 놓친 것 한 가지와 그 이유 \src{(예: D-3의 원금 보장, B의 숫자)}
\item 2일차에 확인하고 싶은 질문
\end{enumerate}
\vspace{0.6em}
\begin{alertblock}{2일차에는 Azure를 쓰지 않는다 — 수업 후 지워도 된다}
VS Code 터미널에서 \texttt{az group delete -n jpm-course -{}-yes} \src{(내 DB와 Azure OpenAI가 함께 지워진다)}
\end{alertblock}
\end{frame}
